% Options for packages loaded elsewhere
\PassOptionsToPackage{unicode}{hyperref}
\PassOptionsToPackage{hyphens}{url}
\documentclass[
  english,
]{article}
\usepackage{xcolor}
\usepackage{amsmath,amssymb}
\setcounter{secnumdepth}{-\maxdimen} % remove section numbering
\usepackage{iftex}
\ifPDFTeX
  \usepackage[T1]{fontenc}
  \usepackage[utf8]{inputenc}
  \usepackage{textcomp} % provide euro and other symbols
\else % if luatex or xetex
  \usepackage{unicode-math} % this also loads fontspec
  \defaultfontfeatures{Scale=MatchLowercase}
  \defaultfontfeatures[\rmfamily]{Ligatures=TeX,Scale=1}
\fi
\usepackage{lmodern}
\ifPDFTeX\else
  % xetex/luatex font selection
\fi
% Use upquote if available, for straight quotes in verbatim environments
\IfFileExists{upquote.sty}{\usepackage{upquote}}{}
\IfFileExists{microtype.sty}{% use microtype if available
  \usepackage[]{microtype}
  \UseMicrotypeSet[protrusion]{basicmath} % disable protrusion for tt fonts
}{}
\makeatletter
\@ifundefined{KOMAClassName}{% if non-KOMA class
  \IfFileExists{parskip.sty}{%
    \usepackage{parskip}
  }{% else
    \setlength{\parindent}{0pt}
    \setlength{\parskip}{6pt plus 2pt minus 1pt}}
}{% if KOMA class
  \KOMAoptions{parskip=half}}
\makeatother
\usepackage{longtable,booktabs,array}
\usepackage{caption}
\captionsetup[table]{skip=6pt}
\newcounter{none} % for unnumbered tables
\usepackage{calc} % for calculating minipage widths
% Correct order of tables after \paragraph or \subparagraph
\usepackage{etoolbox}
\makeatletter
\patchcmd\longtable{\par}{\if@noskipsec\mbox{}\fi\par}{}{}
\makeatother
% Allow footnotes in longtable head/foot
\IfFileExists{footnotehyper.sty}{\usepackage{footnotehyper}}{\usepackage{footnote}}
\makesavenoteenv{longtable}
% definitions for citeproc citations
\NewDocumentCommand\citeproctext{}{}
\NewDocumentCommand\citeproc{mm}{%
  \begingroup\def\citeproctext{#2}\cite{#1}\endgroup}
\makeatletter
 % allow citations to break across lines
 \let\@cite@ofmt\@firstofone
 % avoid brackets around text for \cite:
 \def\@biblabel#1{}
 \def\@cite#1#2{{#1\if@tempswa , #2\fi}}
\makeatother
\newlength{\cslhangindent}
\setlength{\cslhangindent}{1.5em}
\newlength{\csllabelwidth}
\setlength{\csllabelwidth}{3em}
\newenvironment{CSLReferences}[2] % #1 hanging-indent, #2 entry-spacing
 {\begin{list}{}{%
  \setlength{\itemindent}{0pt}
  \setlength{\leftmargin}{0pt}
  \setlength{\parsep}{0pt}
  % turn on hanging indent if param 1 is 1
  \ifodd #1
   \setlength{\leftmargin}{\cslhangindent}
   \setlength{\itemindent}{-1\cslhangindent}
  \fi
  % set entry spacing
  \setlength{\itemsep}{#2\baselineskip}}}
 {\end{list}}
\usepackage{calc}
\newcommand{\CSLBlock}[1]{\hfill\break\parbox[t]{\linewidth}{\strut\ignorespaces#1\strut}}
\newcommand{\CSLLeftMargin}[1]{\parbox[t]{\csllabelwidth}{\strut#1\strut}}
\newcommand{\CSLRightInline}[1]{\parbox[t]{\linewidth - \csllabelwidth}{\strut#1\strut}}
\newcommand{\CSLIndent}[1]{\hspace{\cslhangindent}#1}
\ifLuaTeX
\usepackage[bidi=basic,shorthands=off]{babel}
\else
\usepackage[bidi=default,shorthands=off]{babel}
\fi
\ifLuaTeX
  \usepackage{selnolig} % disable illegal ligatures
\fi
\setlength{\emergencystretch}{3em} % prevent overfull lines
\providecommand{\tightlist}{%
  \setlength{\itemsep}{0pt}\setlength{\parskip}{0pt}}

\usepackage{booktabs}
\usepackage{amsmath,amssymb}
\usepackage{geometry}
\geometry{a4paper, margin=1in}
\usepackage{microtype}
\usepackage{fancyhdr}
\usepackage{xcolor}

\definecolor{academicblue}{RGB}{20, 50, 95}

\usepackage{hyperref}
\hypersetup{
    colorlinks=true,
    linkcolor=academicblue,
    citecolor=academicblue,
    urlcolor=academicblue
}

\pagestyle{fancy}
\fancyhf{}
\renewcommand{\headrulewidth}{0.4pt}
\rhead{\small\textcolor{gray}{\leftmark}}
\cfoot{\thepage}
\usepackage{bookmark}
\IfFileExists{xurl.sty}{\usepackage{xurl}}{} % add URL line breaks if available
\urlstyle{same}
% fallback for those not using the hyperref driver hyperxmp:
\makeatletter
\@ifundefined{xmpquote}{\newcommand{\xmpquote}[1]{#1}}{}
\makeatother
\hypersetup{
  pdftitle={Polymorphic Decision Protocol: A Normative Procedure for High-Uncertainty Decisions with AI Panels},
  pdfauthor={Zachary Shallbetter},
  pdflang={en},
  pdfkeywords={\xmpquote{decision protocols}, \xmpquote{language
models}, \xmpquote{forecast
combination}, \xmpquote{calibration}, \xmpquote{human
oversight}, \xmpquote{research methods}},
  hidelinks,
  pdfcreator={LaTeX via pandoc}}

\title{Polymorphic Decision Protocol: A Normative Procedure for
High-Uncertainty Decisions with AI Panels}
\author{Zachary Shallbetter}
\date{2026-11-01}

\begin{document}
\maketitle
\begin{abstract}
This Type II protocol paper analyzes the normative design, assumptions,
implementation boundary, and falsifiability of PDP v1.3.0. It reports no
efficacy result. A private author-reported case is included only as
development context and is not independently verified. Test 002 is
unexecuted.
\end{abstract}

{
\setcounter{tocdepth}{3}
\tableofcontents
}
\section{00. Implementation status
matrix}\label{implementation-status-matrix}

{\def\LTcaptype{none} % do not increment counter
\begin{longtable}[]{@{}
  >{\raggedright\arraybackslash}p{(\linewidth - 4\tabcolsep) * \real{0.3333}}
  >{\raggedright\arraybackslash}p{(\linewidth - 4\tabcolsep) * \real{0.3333}}
  >{\raggedright\arraybackslash}p{(\linewidth - 4\tabcolsep) * \real{0.3333}}@{}}
\toprule\noalign{}
\begin{minipage}[b]{\linewidth}\raggedright
Capability
\end{minipage} & \begin{minipage}[b]{\linewidth}\raggedright
Status at frozen repository baseline
\end{minipage} & \begin{minipage}[b]{\linewidth}\raggedright
What the status establishes
\end{minipage} \\
\midrule\noalign{}
\endhead
\bottomrule\noalign{}
\endlastfoot
Normative decision procedure & Documented at protocol v1.3.0 &
Requirements exist in prose; no conformance determination is implied \\
Mathematical routines & Implemented and unit-tested in Rust & Only
tested code paths are covered; no empirical decision validity follows \\
UI behavior & Implemented and tested in TypeScript & Local interface
behavior only \\
Multi-provider orchestration & Design only & Not represented as
implemented \\
Durable weight custody and append-only audit & Proposed schema and
architecture & No deployed access-control or operational evidence
supplied \\
Pressure-response classification & Normative proposal & No validated
instrument, reliability estimate, or cross-vendor corpus supplied \\
Prospective performance comparison & Planned, not executed & Test 002
has no results \\
\end{longtable}
}

The matrix separates \textbf{documented}, \textbf{implemented},
\textbf{tested}, and \textbf{empirically validated}. These states are
not interchangeable. The current paper reports no empirical validation.

\section{1. Abstract and executive
decision}\label{abstract-and-executive-decision}

\subsection{Abstract}\label{abstract}

The Polymorphic Decision Protocol (PDP) is a normative procedure for
organizing high-uncertainty decisions made with AI panels while
reserving decision authority for a human. This paper examines the
protocol as a designed method, not as an evaluated intervention. We
reconstruct its four decision phases and portfolio-governance extension,
identify the assumptions behind its consensus and pressure-response
rules, distinguish documented requirements from the limited software
implementation, and specify empirical conditions under which the
method's claims could fail. The current protocol proposes
provenance-aware panel formation, ordered adversarial review, frozen
human-set criteria, commitment and invalidation rules, and outcome-blind
audit. Its adjusted convergence score uses fixed categorical weights
based on vendor and evidence-class differences; those weights are not
estimates of error dependence. The protocol's effectiveness, measurement
validity, safety, and generality remain unestablished. A private
author-reported World Cup case is mentioned only as development context
and is not independently verified. The planned Test 002 has not been
executed. The paper's defensible contribution is therefore a bounded
protocol specification and test agenda. It makes no claim that PDP
improves forecast accuracy, prevents sycophancy, or outperforms a
baseline.

\textbf{Keywords:} decision protocols; language models; forecast
combination; calibration; human oversight; research methods

\subsection{Executive decision}\label{executive-decision}

PDP v1.3.0 is a normative candidate that specifies a potentially
testable sequence of decision controls. The available record supports
describing what the specification requires and what parts of the
repository implement. It does not support a conclusion that the
procedure improves decisions, calibrates AI systems, identifies
sycophancy reliably, or transfers across domains. The immediate research
decision is to submit the protocol and its measurement assumptions to
independent expert review, then freeze a study-specific protocol before
any prospective evaluation.

\subsection{Author and status}\label{author-and-status}

Zachary Shallbetter, Independent Researcher; ORCID:
\href{https://orcid.org/0009-0009-9450-3429}{0009-0009-9450-3429}.
Manuscript version 1.0.0; protocol baseline 1.3.0. Public pre-review
preprint released 2026-11-01
(\href{https://github.com/zachshallbetter/polymorphic-decision-protocol/releases/tag/v1.4.1}{release
page}). It has not received independent review or peer review. No DOI is
assigned. Author verification of claims, mathematics,
implementation-status descriptions, and citations remains outstanding.

\subsection{Claim boundary}\label{claim-boundary}

The paper makes three bounded claims: (1) the frozen specification
documents a staged procedure; (2) the repository contains tested
components but not an end-to-end implementation; and (3) comparative
efficacy is unresolved. These are claims about documents, code, and
missing admissible evidence. They are not claims about causal effects or
real-world performance.

\section{2. Introduction}\label{introduction}

\subsection{2.1 Research question}\label{research-question}

The immediate question is descriptive and methodological: what decision
procedure does PDP v1.3.0 specify, what assumptions do its controls
require, and what evidence would be necessary to evaluate them? The
future empirical question is whether, under a defined decision
population and comparison design, use of PDP changes decision quality or
calibration relative to a prespecified baseline. This paper answers only
the first question.

\subsection{2.2 Contribution type}\label{contribution-type}

The paper is a Type II protocol specification. It describes a proposed
method, maps its normative rules to an implementation boundary, and
identifies measurable failure conditions. It does not report a completed
study, a validated instrument, a causal model, or a validated
intervention. Its contribution claim is limited to presenting and
critically bounding the existing v1.3.0 protocol in a scholarly form.

\subsection{2.3 Decision standard}\label{decision-standard}

Protocol review should establish whether requirements are internally
coherent, measurable, feasible, and ethically acceptable. Empirical
evaluation requires a separate preregistration that specifies
population, sampling, treatment, comparator, outcomes, estimand,
stopping rule, uncertainty procedure, and missing-data handling. Neither
successful software tests nor expert agreement can substitute for that
evaluation.

\subsection{2.4 Core explanandum}\label{core-explanandum}

The core explanandum is the design of a decision procedure that
constrains how panel outputs are elicited, challenged, and recorded
while retaining human authority over value choices and commitment. This
paper describes that proposed procedure; it does not explain or measure
real-world decision outcomes.

\section{2.5 Scope and exclusions}\label{scope-and-exclusions}

\subsection{Scope}\label{scope}

PDP concerns decisions with uncertain outcomes and a commitment point.
Its normative text organizes model elicitation, evidence review,
pressure testing, human custody of weights and stakes, commitment
controls, monitoring, and post-outcome audit. The protocol can be
analyzed as a governance procedure even when it is not implemented as
software.

\subsection{Exclusions}\label{exclusions}

This paper does not establish or claim: predictive accuracy; causal
improvement over single-model or human baselines; prevalence of
sycophancy; validity of the adjusted convergence score; safety for
financial, medical, legal, hiring, or other high-impact use; transfer
across domains; correctness of private source data; or public validation
of the originating case. The current code is not an autonomous trading
system. No live execution or recommendation should be inferred from this
paper.

\subsection{Development illustration}\label{development-illustration}

The protocol originated in part from an author-reported 2026 World Cup
decision involving a multi-model panel. The underlying conversations and
trade records remain private and are not included here. The account is
not independently inspectable, was not collected under the proposed
protocol, and is vulnerable to recall, selection, and outcome bias. It
is cited only as a development illustration. No numerical case result is
reproduced, and no claim about the case is treated as verified evidence.
Test 002 remains planned and unexecuted.

\section{3. Prior art and explanatory
boundaries}\label{prior-art-and-explanatory-boundaries}

Probabilistic scoring has a mature foundation. Brier's score evaluates
the accuracy of probability forecasts against realized outcomes (Brier
1950). Forecast combination also predates language models; Bates and
Granger showed that combinations can reduce mean-square error in a
particular forecasting context, with weights informed by historical
forecast errors (Bates and Granger 1969). PDP therefore cannot claim
novelty for probability scoring, forecast combination, or the general
insight that dependence matters.

Research on multi-agent debate evaluates how interaction strategies
affect model performance, with results varying by task and strategy
(Smit et al. 2024). Research on sycophancy studies model behavior that
aligns with user views at the expense of truthfulness under specified
tasks ({Sharma et al.} 2023). These literatures make the motivating
concerns scientifically plausible research topics. They do not establish
that PDP's exact prompt, categorical flip labels, or sequence controls
measure the same constructs or improve outcomes.

PDP combines governance rules and operational proposals: evidence-class
restrictions, a two-axis categorical agreement adjustment, a challenge
prompt, frozen human-set weights, commitment limits, and audit
procedures. This paper does not claim that their combination is novel or
more effective. A broader systematic prior-art search and independent
novelty assessment remain required. The source register identifies the
records checked for the cited bibliographic details; those records do
not establish PDP's empirical effectiveness.

\section{4. Method and evidential
basis}\label{method-and-evidential-basis}

\subsection{4.1 Document analysis}\label{document-analysis}

This paper treats the v1.3.0 normative specification as its object
(Shallbetter 2026). The baseline is repository commit
\texttt{35d945a04e622a560c8e1be53ffc0a796be1aa79}, corresponding to the
source snapshot associated with repository release v1.3.2. The
manuscript itself is a later analysis and does not alter protocol rules.
Repository code and tests are inspected as implementation artifacts, not
as evidence of empirical effectiveness.

\subsection{4.2 Evidence lanes}\label{evidence-lanes}

Statements about protocol requirements are document-grounded. Statements
about software are repository-grounded and bounded to inspected source
and tests. Claims about effectiveness, prevalence, user outcomes, or
generality are pending proof. The private originating case is not
admissible for independent verification in this paper. No new experiment
or participant data were collected for this manuscript.

\subsection{4.3 Freeze and registration}\label{freeze-and-registration}

The normative baseline is frozen for this analysis at the commit above.
The paper's claims and Test 002 prediction are recorded in the RDP
registry, but the present working tree is uncommitted. The prediction is
therefore explicitly marked pending registration. It must be committed
before Test 002 starts; there is no claim of commit-ordered
preregistration yet.

\subsection{4.4 Review method}\label{review-method}

Two independent reviews are planned: one on research design and
measurement, and one on AI evaluation or protocol methods. Reviewers
have not been contacted and no feedback has been received. The review
packet is prepared separately. Reviewer selection, independence,
conflicts, and dispositions must be recorded before any
reviewed-preprint status is used. The public GitHub version is a
pre-review preprint, not peer reviewed; its release does not imply
independent review.

\section{4.5 Minimum formal model}\label{minimum-formal-model}

Let decision instance \(d\) have candidate set \(A_d\), declared
information state \(I_d\), human criteria weights \(w_d\), panel agents
\(K_d\), and eventual outcome \(y_d\). Agent \(k\) returns forecast
distribution \(p_{kd}\) after observing its permitted evidence
\(E_{kd} \subseteq I_d\). The protocol records a decision procedure
\(\pi\) mapping the declared inputs and responses to a recommendation
record, commitment gate, and audit record.

For a categorical outcome with \(J\) possible values, the multicategory
Brier score is

\[BS(p,y)=\sum_{j=1}^{J}(p_j-\mathbb{1}[y=j])^2.\]

For a binary event this can be reported in the convention
\(BS=(p-y)^2\). Lower scores are better for the observed forecast set,
but a sample estimate requires a defined sampling process and
uncertainty analysis. A skill score against baseline forecasts \(q_d\)
may be written

\[BSS=1-\frac{\sum_d BS(p_d,y_d)}{\sum_d BS(q_d,y_d)}.\]

It is undefined when the baseline denominator is zero and can be
unstable in small or dependent samples. This paper reports no Brier
scores and no Brier skill score.

The v1.3.0 adjusted convergence score is a categorical rule over
agreement indicators and declared weights. Its weights are not
identified from a covariance model, and the current paper does not
reinterpret them as such. No theorem in this paper establishes that the
score estimates consensus reliability or decision quality.

\section{5. Proposed evaluation and
falsifiability}\label{proposed-evaluation-and-falsifiability}

\subsection{5.1 Decisive tests}\label{decisive-tests}

{\def\LTcaptype{none} % do not increment counter
\begin{longtable}[]{@{}
  >{\raggedright\arraybackslash}p{(\linewidth - 4\tabcolsep) * \real{0.3333}}
  >{\raggedright\arraybackslash}p{(\linewidth - 4\tabcolsep) * \real{0.3333}}
  >{\raggedright\arraybackslash}p{(\linewidth - 4\tabcolsep) * \real{0.3333}}@{}}
\toprule\noalign{}
\begin{minipage}[b]{\linewidth}\raggedright
Claim under test
\end{minipage} & \begin{minipage}[b]{\linewidth}\raggedright
Design needed
\end{minipage} & \begin{minipage}[b]{\linewidth}\raggedright
Failure condition
\end{minipage} \\
\midrule\noalign{}
\endhead
\bottomrule\noalign{}
\endlastfoot
Pressure-response labels are reproducible & Blinded raters, controlled
repeated prompts, frozen context, prespecified coding manual &
Reliability is inadequate or labels change materially under irrelevant
wording changes \\
Categorical agreement adjustment tracks dependence or improves decisions
& Repeated forecasts, outcome data, comparison with explicit dependence
estimates and simple baselines & No incremental predictive or decision
value, or weights worsen calibration under the prespecified analysis \\
Procedure improves forecast quality & Prospective sample, registered
estimand, comparator, complete outcomes, uncertainty intervals & PDP is
not better than the baseline under the declared decision criterion, or
causes material harms \\
Human custody controls are operationally followed & Independent audit of
frozen criteria, overrides, and commitment records & Unlogged changes
occur or the procedure cannot preserve the stated authority boundary \\
\end{longtable}
}

These are proposed tests, not findings. Sample size, population,
outcomes, and statistical decision rules must be specified before
execution. A single case cannot establish broad performance.

\subsection{5.2 Limitations}\label{limitations}

The protocol is a normative candidate. The literature search is
preliminary. The private originating case cannot be audited here. No
prospective study, independent software review, or protocol conformance
assessment is reported. The categorical consensus weights lack empirical
calibration. The challenge-based classifier may confound multiple causes
of response change. Existing tests establish only behavior of tested
code paths. No operational safety or regulatory suitability claim is
made.

\subsection{5.3 Author and funding
disclosures}\label{author-and-funding-disclosures}

The sole author originated PDP and is not independent of its design.
Author contributions under the CRediT taxonomy, funding status, and any
additional competing-interest declarations remain unconfirmed; no
unverified declaration is made here. The author is identified as an
Independent Researcher. This public pre-review preprint was released
before those confirmations; the release does not imply that the
declarations are complete.

\subsection{5.4 Data, code, and AI
assistance}\label{data-code-and-ai-assistance}

No new study data were collected. Private originating-case records and
transcripts are not included and are unavailable for independent
verification. PDP source code is available in the associated repository
under its software license; this paper is licensed separately under CC
BY 4.0. OpenAI Codex assisted with repository inspection, literature
discovery, critical analysis, manuscript editing, and preparation of
research records and release artifacts. It is not an author. Author
verification of claims, mathematics, code, and citations remains
outstanding in this pre-review version.

\subsection{5.5 Review and publication
status}\label{review-and-publication-status}

Two independent expert reviews are planned and have not yet been
obtained. This is a public pre-review preprint; it is not independently
reviewed or peer reviewed. Reviewer comments will not be represented as
empirical validation.

\section{6. Results and negative
evidence}\label{results-and-negative-evidence}

This analysis generated no empirical results. No participants, model
outputs, forecasts, or outcomes were collected or reanalyzed. The
author-reported development case is not independently inspectable and is
not used as a result.

The repository inspection identified a difference between the normative
scope and the implementation: mathematical routines and UI behavior have
code and tests, while multi-provider orchestration, durable governance
controls, and end-to-end audit remain design proposals. This is a status
finding about repository contents, not a performance finding.

No admissible positive or negative efficacy result is available. The
research record therefore treats the central performance question as
unresolved and defines conditions that could falsify the proposed
measurement and outcome claims.

\section{7. Discussion: protocol analysis, ontology, and
assumptions}\label{discussion-protocol-analysis-ontology-and-assumptions}

\subsection{Stages}\label{stages}

The v1.3.0 specification describes thesis formation, adversarial stress
testing, structured commitment, monitoring and settlement, and portfolio
governance across concurrent decisions. Stage gates are intended to
prevent commitment when a required step fails. The paper reports these
as normative design features, not observed operating properties.

\subsection{Human authority}\label{human-authority}

The specification reserves value weights, stakes, and exit decisions to
a human. This is a governance allocation. It does not establish that
human judgments are accurate, unbiased, or legally sufficient. A frozen
weight vector can preserve an invalid value model just as faithfully as
a defensible one; the specification must therefore be evaluated
separately for procedural integrity and decision quality.

\subsection{Pressure response}\label{pressure-response}

The protocol proposes an ordered challenge without new evidence and
labels response patterns such as robust, evidence-driven flip,
pressure-driven flip, and collapse. Such labels require an operational
coding rule. The presence of a challenge alone cannot identify the cause
of an answer change: stochastic variation, context effects, new facts,
or ambiguity may also matter. The classifier needs blinded human or
validated automated scoring, reliability estimates, and controls before
it can be treated as a measurement instrument.

\subsection{Consensus adjustment}\label{consensus-adjustment}

The current specification assigns categorical weights according to
vendor and evidence-class differences. The values are design choices,
not empirically estimated dependence coefficients. Distinct vendors do
not guarantee independent errors, and shared evidence does not fully
determine error correlation. Until calibrated against outcome data, the
score should be interpreted as a declared governance heuristic, not as a
probability, effective sample size, or validated measure of epistemic
independence.

\subsection{Commitment and audit}\label{commitment-and-audit}

Precommitted conditions and outcome-blind audit can make decision
records more inspectable. They also introduce operational burdens and
risks: ambiguous triggers, missing records, changing environments, and
pressure to treat a mechanical threshold as objective authority. The
protocol should report compliance separately from accuracy and utility.

\subsection{Canonical ontology}\label{canonical-ontology}

{\def\LTcaptype{none} % do not increment counter
\begin{longtable}[]{@{}
  >{\raggedright\arraybackslash}p{(\linewidth - 6\tabcolsep) * \real{0.2500}}
  >{\raggedright\arraybackslash}p{(\linewidth - 6\tabcolsep) * \real{0.2500}}
  >{\raggedright\arraybackslash}p{(\linewidth - 6\tabcolsep) * \real{0.2500}}
  >{\raggedright\arraybackslash}p{(\linewidth - 6\tabcolsep) * \real{0.2500}}@{}}
\toprule\noalign{}
\begin{minipage}[b]{\linewidth}\raggedright
Term
\end{minipage} & \begin{minipage}[b]{\linewidth}\raggedright
Type
\end{minipage} & \begin{minipage}[b]{\linewidth}\raggedright
Definition in this analysis
\end{minipage} & \begin{minipage}[b]{\linewidth}\raggedright
Exclusion
\end{minipage} \\
\midrule\noalign{}
\endhead
\bottomrule\noalign{}
\endlastfoot
Candidate decision & Decision object & A bounded choice under
uncertainty before commitment & Not an AI recommendation or realized
outcome \\
Claim & Proposition & A statement whose support and failure condition
can be assessed & Not a preference or confidence expression alone \\
Evidence class & Input constraint & A declared category restricting what
an agent may use & Not proof that two agents' errors are independent \\
Agreement adjustment & Governance heuristic & The v1.3.0 categorical
weighting of agreement based on vendor and evidence-class differences &
Not a calibrated probability or dependence estimate \\
Pressure response & Observed response category & A coded change or
stability after a declared challenge & Not proof of sycophancy without
validated design \\
Frozen weights & Human-set vector & Criteria weights recorded before
evaluation and not silently changed & Not evidence that the weights are
correct \\
Invalidation condition & Prospective rule & An observable condition
paired with a predeclared response & Not an agent utterance or
retrospective explanation \\
Protocol compliance & Process measure & Whether declared procedure
requirements were followed & Not decision accuracy, safety, or
utility \\
\end{longtable}
}

\subsection{Assumptions and first
principles}\label{assumptions-and-first-principles}

PDP assumes that input provenance can be described well enough to inform
panel construction; criteria and weights can be set before model scores
are observed; a challenge can be applied without adding evidence;
response categories can be distinguished reliably; commitment and
invalidation can be operationalized; and outcomes can be scored on a
declared horizon. Each can fail in practice and requires its own
evidence.

The human-authority principle is normative, not an empirical axiom. It
assigns authority but does not establish that the decision owner has
expertise, independence, or adequate information. A frozen weight vector
can preserve an invalid value model as faithfully as a defensible one;
procedural integrity and decision quality must therefore be evaluated
separately.

The model-independence assumption is particularly uncertain. Provider
and evidence labels are observable proxies, while residual dependence
must be measured from repeated forecasts and outcomes. The current
protocol supplies no validated dependence model.

\section{8. Conclusion}\label{conclusion}

PDP v1.3.0 specifies a staged governance procedure for uncertain
decisions involving AI panels. Its strongest present contribution is a
set of explicit procedural requirements that can be inspected,
challenged, and converted into future tests. The present record does not
show that the procedure improves accuracy, calibration, safety, or
utility. The categorical agreement adjustment and pressure-response
labels require measurement validation, while code tests must remain
distinct from system-level and empirical evidence.

The next defensible step is critical review of the protocol's research
design and measurement assumptions. Any empirical study must be frozen
before execution and must report outcomes irrespective of direction.
Until that work is complete, PDP should be described as a normative
candidate and the originating case as an unverified author report.

\section{9. Illustrative example}\label{illustrative-example}

\subsection{9.1 Author-reported
illustration}\label{author-reported-illustration}

An author-reported World Cup decision motivated several protocol rules.
The records are private and unavailable for independent inspection. The
decision was not prospectively conducted under the final protocol, and
the author both designed the rules and interprets the case. It is
therefore a development illustration only. No outcome, score, or quoted
transcript is presented as verified evidence. The case cannot establish
that any rule caused a decision, detected a model behavior, or
generalizes.

\subsection{9.2 Implementation-status
matrix}\label{implementation-status-matrix-1}

{\def\LTcaptype{none} % do not increment counter
\begin{longtable}[]{@{}
  >{\raggedright\arraybackslash}p{(\linewidth - 4\tabcolsep) * \real{0.3333}}
  >{\raggedright\arraybackslash}p{(\linewidth - 4\tabcolsep) * \real{0.3333}}
  >{\raggedright\arraybackslash}p{(\linewidth - 4\tabcolsep) * \real{0.3333}}@{}}
\toprule\noalign{}
\begin{minipage}[b]{\linewidth}\raggedright
Capability
\end{minipage} & \begin{minipage}[b]{\linewidth}\raggedright
Repository status
\end{minipage} & \begin{minipage}[b]{\linewidth}\raggedright
Evidence boundary
\end{minipage} \\
\midrule\noalign{}
\endhead
\bottomrule\noalign{}
\endlastfoot
Core mathematical routines & Implemented in Rust & Unit tests exercise
the implemented functions; no empirical validity follows \\
UI behavior & Implemented in TypeScript & UI tests cover local behavior
only \\
Multi-provider orchestration & Design document only & No end-to-end
orchestration is claimed \\
Durable custody and append-only audit controls & Proposed schema/design
& Schema presence is not proof of deployed permissions or operational
control \\
Pressure-response classification across vendors & Proposed method & No
validated corpus, inter-rater study, or cross-vendor evaluation
supplied \\
Prospective performance comparison & Planned, not run & Test 002 remains
unexecuted \\
\end{longtable}
}

\subsection{9.3 Negative evidence}\label{negative-evidence}

No admissible negative or positive empirical result is reported. The
absence of admissible evidence is not evidence that the protocol
succeeds or fails. The negative-results register records the open
prediction and the conditions that would expose measurement failure.

\section{10. Concept and implementation audit, limitations, and
disclosures}\label{concept-and-implementation-audit-limitations-and-disclosures}

{\def\LTcaptype{none} % do not increment counter
\begin{longtable}[]{@{}
  >{\raggedright\arraybackslash}p{(\linewidth - 6\tabcolsep) * \real{0.2500}}
  >{\raggedright\arraybackslash}p{(\linewidth - 6\tabcolsep) * \real{0.2500}}
  >{\raggedright\arraybackslash}p{(\linewidth - 6\tabcolsep) * \real{0.2500}}
  >{\raggedright\arraybackslash}p{(\linewidth - 6\tabcolsep) * \real{0.2500}}@{}}
\toprule\noalign{}
\begin{minipage}[b]{\linewidth}\raggedright
Protocol concept
\end{minipage} & \begin{minipage}[b]{\linewidth}\raggedright
Normative location
\end{minipage} & \begin{minipage}[b]{\linewidth}\raggedright
Software evidence
\end{minipage} & \begin{minipage}[b]{\linewidth}\raggedright
Unresolved issue
\end{minipage} \\
\midrule\noalign{}
\endhead
\bottomrule\noalign{}
\endlastfoot
Adjusted convergence score & PDP v1.3.0, panel formation & Mathematical
subset in \texttt{pdp-core}; unit tests & Fixed category weights are not
estimates of error dependence \\
Pressure-response classes & PDP v1.3.0, stress test & No validated
cross-vendor classifier & Coding reliability, stochastic response, and
alternative causes \\
Frozen human weights & PDP v1.3.0, formal model and commitment &
Repository code does not constitute a complete custody service & No
deployed system or independent control audit \\
Invalidation and commitment gates & PDP v1.3.0, commitment and
monitoring & Architecture/schema proposals & Ambiguity, authority, and
real-world enforceability \\
Calibration audit & PDP v1.3.0, outcome audit & Brier-related math is
not a longitudinal study & Baseline definition, sample dependence,
decision utility \\
Portfolio governance & PDP v1.3.0, concurrent decisions & No complete
service implementation & Exposure correlation and operational
estimation \\
\end{longtable}
}

The code and architecture are not interchangeable. Where the
architecture describes components absent from code, this paper calls
them design proposals. The paper does not treat a schema as an enforced
permission boundary unless deployed database grants and operational
tests demonstrate that behavior.

\subsection{Disclosures and publication
conditions}\label{disclosures-and-publication-conditions}

The sole author originated PDP and is not independent of its design. The
author is identified as Zachary Shallbetter, Independent Researcher
(ORCID 0009-0009-9450-3429). \textbf{CRediT:} role assignments have not
been confirmed by the author; no roles are assigned. \textbf{Funding:}
status requires author confirmation; no funding declaration is made.
\textbf{Competing interests:} the author developed PDP and has an
intellectual interest in its evaluation; any additional financial or
nonfinancial interests remain unconfirmed.

No new study data were collected. Private originating-case records and
transcripts are not included and are unavailable for independent
verification. PDP source code is available at
\href{https://github.com/zachshallbetter/polymorphic-decision-protocol}{the
public repository} under its MIT license. This paper is licensed under
CC BY 4.0. OpenAI Codex assisted with repository inspection, literature
discovery, critical analysis, manuscript editing, and preparation of
research records and release artifacts. It is not an author. Author
verification of claims, mathematics, code, and citations remains
outstanding in this public pre-review version.

Two independent expert reviews are planned and have not yet been
obtained. This is a public pre-review preprint; it is not independently
reviewed or peer reviewed. Reviewer comments will not be represented as
empirical validation.

\subsection{References}\label{references}

The bibliography is provided in the accompanying \texttt{references.bib}
file and is rendered from citation keys in the Markdown and PDF sources.

\protect\phantomsection\label{refs}
\begin{CSLReferences}{1}{1}
\bibitem[\citeproctext]{ref-batesgranger1969}
Bates, J. M., and C. W. J. Granger. 1969. {``The Combination of
Forecasts.''} \emph{Operational Research Quarterly} 20 (4): 451--68.
\url{https://doi.org/10.1057/jors.1969.103}.

\bibitem[\citeproctext]{ref-brier1950}
Brier, Glenn W. 1950. {``Verification of Forecasts Expressed in Terms of
Probability.''} \emph{Monthly Weather Review} 78 (1): 1--3.
\url{https://doi.org/10.1175/1520-0493(1950)078\%3C0001:VOFEIT\%3E2.0.CO;2}.

\bibitem[\citeproctext]{ref-pdpv130}
Shallbetter, Zachary. 2026. \emph{Polymorphic Decision Protocol}. V.
1.3.2. Zenodo, released. \url{https://doi.org/10.5281/zenodo.23798205}.

\bibitem[\citeproctext]{ref-sharma2023}
{Sharma, Mrinank, Meg Tong, Tomasz Korbak, et al.} 2023. \emph{Towards
Understanding Sycophancy in Language Models}.
\url{https://arxiv.org/abs/2310.13548}.

\bibitem[\citeproctext]{ref-smit2024}
Smit, Andries Petrus, Nathan Grinsztajn, Paul Duckworth, Thomas D.
Barrett, and Arnu Pretorius. 2024. {``Should We Be Going MAD? A Look at
Multi-Agent Debate Strategies for LLMs.''} \emph{Proceedings of the 41st
International Conference on Machine Learning}, Proceedings of machine
learning research, vol. 235: 45883--905.
\url{https://proceedings.mlr.press/v235/smit24a.html}.

\end{CSLReferences}

\end{document}
